This section talks about the process of the RQL reaction mechanism.





Boundary conditions:




T in fuel   = 500 K

T in air    = 800 K  

P in fuel   = 15 bar

P in air    = 12 bar  




%\begin{figure}[H]
%    \centering
%    \includesvg[width=0.7\linewidth]{figs/rql_process_diagram.svg}
%    \caption{Caption}
%    \label{fig:enter-label}
%\end{figure}


Writing the balanced chemical equation

\begin{equation}
    \ce{4NH3 + 3(O2 + 3.76 N2) -> 2N2 + 6H2O + 11.28N2}\label{eq:balanced_reaction}
\end{equation}

In \autoref{eq:balanced_reaction}, the oxidizer nitrogen is kept separate from the inert air nitrogen to keep track.
